% LNCS-style paper: Trust-Aware Query Optimization for AI-Native Data Systems
% Compile with: pdflatex trust_aware_paper.tex (requires llncs.cls)

\documentclass[runningheads]{llncs}

\usepackage[T1]{fontenc}
\usepackage{graphicx}
\usepackage{amsmath}
\usepackage{amssymb}
\usepackage{booktabs}
\usepackage{hyperref}
\usepackage{microtype}
\usepackage{xcolor}
\usepackage{listings}
\usepackage{algorithm}
\usepackage{algorithmic}

\lstset{
  basicstyle=\ttfamily\footnotesize,
  breaklines=true,
  frame=single,
  captionpos=b,
}

\begin{document}

\title{Trust-Aware Query Optimization\\for AI-Native Data Systems}

\titlerunning{Trust-Aware Query Optimization}

\author{Jorge Martinez-Gil\inst{1}}
\authorrunning{J. Martinez-Gil}

\institute{
  Software Competence Center Hagenberg GmbH (SCCH), Austria\\
  \email{jorge.martinez-gil@scch.at}
}

\maketitle

%------------------------------------------------------------------------
\begin{abstract}
Modern AI-native data systems expose a heterogeneous mix of query
operators---relational tables, vector indexes, retrieval services,
model-backed tools, memory shards, and grounded generation units---that
differ not only in cost and latency but also in trustworthiness,
calibration confidence, privacy exposure, and hallucination risk.
Classical cost-based optimizers were not designed for this landscape:
they minimise execution cost under deterministic quality assumptions that
break down the moment a model call or a retrieval service enters the plan.
We propose \emph{trust-aware query optimization}, a principled extension
to the query planning layer that treats trust as a first-class objective
on a par with cost and latency.
The optimizer enforces hard policy constraints---capability, latency,
freshness, trust floor, privacy ceiling, and hallucination-risk
ceiling---before scoring candidates with a calibrated beta-binomial
trust lower bound combined in a weighted expected-utility function.
Explainable plan objects with Pareto annotations, portable trust
certificates, redundant source portfolios, multi-stage AI workflow
planning, and online feedback loops complete the system.
We evaluate the approach on a deterministic synthetic benchmark with four
ablation policies and on the UCI Wisconsin Diagnostic Breast Cancer
(WDBC) dataset with five-fold cross-validation and Wilson/bootstrap
confidence intervals.
Across both evaluations the trust-aware policy achieves the best
calibrated trust lower bound among utility-maximising policies while
providing the richest audit trail.
All code and data are publicly available as a self-contained,
dependency-free Python artifact.
\end{abstract}

\keywords{query optimization \and trust \and AI-native databases \and
source selection \and explainability \and calibration}

%------------------------------------------------------------------------
\section{Introduction}
\label{sec:intro}

The rise of retrieval-augmented generation (RAG), vector databases, and
LLM-integrated query engines has produced a new class of data
systems~\cite{lewis2020rag,borgeaud2022retro} in which a
\emph{query plan} may route work not just to a relational engine but to
an embedding index, a remote tool API, or a language model call.  Each
such operator carries qualitatively different uncertainty: a table row is
either present or absent, but a model-generated answer may be
confidently wrong.  Existing cost-based optimizers, which date back to
System R~\cite{selinger1979access}, optimize for I/O, CPU, and network
cost under an assumption of deterministic, trustworthy operators.  They
offer no principled mechanism for expressing \emph{``prefer sources
whose past answers were verified by a policy auditor''}  or
\emph{``reject any source whose hallucination risk exceeds 20\%''}.

Trust has been studied in database contexts before: trust management
systems~\cite{blaze1996decentralized} handle authorization, and data
provenance research~\cite{buneman2001provenance} tracks lineage.  Data
quality work~\cite{batini2009methodologies} assigns reliability scores to
sources.  Federated query optimizers~\cite{levy1996querying} select among
heterogeneous sources by schema compatibility and estimated selectivity.
None of these lines of work models \emph{calibrated evidential trust} as
a first-class optimization objective subject to policy-level hard
constraints.

This paper makes the following contributions:
\begin{enumerate}
  \item A data model that captures trust evidence, capabilities,
    freshness, privacy risk, and hallucination risk for each queryable
    source (\S\ref{sec:model}).
  \item A constraint-and-score optimizer that enforces hard policy gates
    before computing calibrated trust lower bounds from beta-binomial
    posteriors (\S\ref{sec:optimizer}).
  \item Named policy profiles that make experiments and deployments
    reproducible (\S\ref{sec:policies}).
  \item Extensions for portfolio planning, multi-stage AI pipeline
    planning, online trust updates, and portable trust certificates
    (\S\ref{sec:extensions}).
  \item Empirical evaluation on a 32-source synthetic benchmark and the
    WDBC real-data study (\S\ref{sec:eval}).
\end{enumerate}

The artifact is intentionally dependency-free.  All statistical
primitives---Wilson score intervals, bootstrap confidence intervals,
beta-binomial posteriors---are implemented from scratch so the code can
run wherever a paper is checked out.

%------------------------------------------------------------------------
\section{Background and Related Work}
\label{sec:related}

\paragraph{Cost-Based Query Optimization.}
The System R optimizer~\cite{selinger1979access} established the
template that all major relational engines follow: enumerate plan
alternatives, estimate their cost using cardinality and operator cost
models, and select the cheapest one.  Extensions handle
parallelism~\cite{mohan1992aries}, approximate
queries~\cite{ioannidis1990randomized}, and
uncertainty~\cite{re2007probabilistic}, but the objective function is
always a form of execution cost, not source trustworthiness.

\paragraph{Trust and Reputation Systems.}
EigenTrust~\cite{kamvar2003eigentrust} and related work model peer
reputation in distributed systems.  Winsborough and Li~\cite{winsborough2002trust}
define trust negotiation protocols.  These systems propagate scalar trust
values but do not integrate trust into relational query planning with hard
policy constraints.

\paragraph{Data Quality and Provenance.}
Batini and Scannapieco~\cite{batini2009methodologies} survey data-quality
dimensions including accuracy, consistency, and timeliness.  Buneman
et al.~\cite{buneman2001provenance} formalize provenance for relational
and semistructured data.  Our \texttt{TrustEvidence} model draws
inspiration from these dimensions but frames them as calibrated evidence
streams amenable to optimization.

\paragraph{Probabilistic and Uncertain Databases.}
MayBMS~\cite{antova2008maybms} and Trio~\cite{widom2005trio} extend the
relational model with tuple-level probabilities.  They model uncertainty
in data, not in the trustworthiness of the operator that produced the
data.

\paragraph{Federated Query Planning.}
Information Manifold~\cite{levy1996querying} selects among heterogeneous
sources by capability and selectivity.  Our work shares the source-selection
framing but adds calibrated trust, policy gates, and Pareto-aware ranking.

\paragraph{LLM Query Engines.}
Recent systems such as LOTUS~\cite{patel2024lotus} and
NL2SQL~\cite{shi2024survey} integrate LLMs into query plans.
They focus on correctness and latency; trust calibration and policy
auditing are left to application code.

%------------------------------------------------------------------------
\section{System Model}
\label{sec:model}

\subsection{Data Sources}

A \emph{data source} $s$ is any queryable unit: a relational table,
vector index, retrieval service, LLM inference endpoint, memory shard,
or grounded generation operator.  Formally, $s$ is characterized by:

\begin{itemize}
  \item $\text{caps}(s)$: a set of capability labels from
    $\{\texttt{structured\_sql},\allowbreak
       \texttt{vector\_search},\allowbreak
       \texttt{lexical\_search},\allowbreak
       \texttt{graph\_lookup},\allowbreak
       \texttt{llm\_inference},\allowbreak
       \texttt{streaming},\allowbreak
       \texttt{freshness\_guarantee},\allowbreak
       \texttt{policy\_audit}\}$.
  \item $\tau(s) \in [0,1]$: a declared trust score.
  \item $\mathcal{E}(s)$: a tuple of \emph{TrustEvidence} items, each
    recording observed successes and failures on a named dimension.
  \item $\ell(s)$: expected latency in milliseconds.
  \item $c(s)$: cost per query in abstract currency units.
  \item $\phi(s) \in [0,1]$: freshness score.
  \item $\pi(s) \in [0,1]$: privacy risk.
  \item $h(s) \in [0,1]$: hallucination risk.
  \item $\rho(s) \in [0,1]$: schema reliability.
\end{itemize}

\subsection{Trust Evidence and Calibration}
\label{sec:calibration}

Each evidence item $e \in \mathcal{E}(s)$ tracks a named dimension
(e.g., \texttt{answer\_acceptance}) with counts $(n^+, n)$ and a
symmetric Beta prior $(\alpha_0, \beta_0)$ defaulting to $(2,2)$.  The
posterior is:
\begin{equation}
  \alpha = \alpha_0 + n^+, \quad
  \beta  = \beta_0  + (n - n^+).
  \label{eq:beta}
\end{equation}
The posterior mean $\mu = \alpha/(\alpha+\beta)$ and variance
$\sigma^2 = \alpha\beta/[(\alpha+\beta)^2(\alpha+\beta+1)]$
give a calibrated lower bound at risk tolerance $\varepsilon$:
\begin{equation}
  \text{lb}(e, z_\varepsilon) =
    \max\!\bigl(0,\; \mu - z_\varepsilon\,\sigma\bigr),
  \label{eq:lb}
\end{equation}
where $z_\varepsilon$ is the one-sided Gaussian quantile for confidence
level $1-\varepsilon$.  The calibrated trust lower bound of a source
combines the declared score with weighted evidence lower bounds:
\begin{equation}
  \text{TLB}(s, \varepsilon) =
    \frac{0.25\,\tau(s) +
          \sum_{e \in \mathcal{E}(s)} w_e\,\text{lb}(e,z_\varepsilon)}
         {0.25 + \sum_{e \in \mathcal{E}(s)} w_e}.
  \label{eq:tlb}
\end{equation}
When no evidence is available the declared score is used unchanged.  The
weight 0.25 anchors the prior weakly: a source with many high-weight
evidence items moves its lower bound toward the observed rate; a source
with no evidence retains its declared score under full uncertainty.

\subsection{Query Requests}

A \emph{query request} $q$ combines hard constraints and soft objective
weights:
\begin{itemize}
  \item Hard: required capabilities $C_q$, maximum latency
    $\ell_{\max}$, minimum trust lower bound $\tau_{\min}$, minimum
    freshness $\phi_{\min}$, maximum cost $c_{\max}$, maximum privacy
    risk $\pi_{\max}$, maximum hallucination risk $h_{\max}$.
  \item Soft: weights
    $(w_\tau, w_\ell, w_\phi, w_c, w_\kappa, w_r, w_u) \geq 0$
    for trust, latency, freshness, cost, coverage, risk, and uncertainty
    penalty respectively.
  \item Risk tolerance $\varepsilon \in [0,1]$ controlling
    $z_\varepsilon$ in Eq.~\ref{eq:lb}.
\end{itemize}

%------------------------------------------------------------------------
\section{Trust-Aware Query Optimizer}
\label{sec:optimizer}

\subsection{Constraint Filtering}

Before any scoring the optimizer rejects all sources that violate a hard
constraint.  Each rejected source records the violated constraint as a
human-readable reason.  This keeps policy boundaries auditable: a cheap
source that violates the privacy ceiling does not win by compensating
with low cost.

\begin{algorithm}[t]
\caption{Trust-Aware Source Selection}
\label{alg:optimizer}
\begin{algorithmic}[1]
\REQUIRE Sources $S$, request $q$
\STATE $\text{admitted} \leftarrow \emptyset$,
       $\text{rejected} \leftarrow \emptyset$
\FOR{each $s \in S$}
  \STATE $R \leftarrow \text{RejectReasons}(s, q)$
  \IF{$R = \emptyset$}
    \STATE $\text{admitted} \leftarrow \text{admitted} \cup \{s\}$
  \ELSE
    \STATE $\text{rejected} \leftarrow \text{rejected} \cup \{(s, R)\}$
  \ENDIF
\ENDFOR
\STATE Score and rank $\text{admitted}$ by Eq.~\ref{eq:score}
\STATE Mark Pareto-optimal candidates
\RETURN explainable \texttt{QueryPlan}
\end{algorithmic}
\end{algorithm}

\subsection{Scoring}

For every admitted source $s$ the optimizer computes a weighted
expected-utility score:
\begin{align}
  \text{score}(s, q) &=
    w_\tau\,\text{TLB}(s,\varepsilon)
    + w_\ell\,\ell\text{-comp}(s,q)
    + w_\phi\,\phi(s) \nonumber\\
  &\quad
    + w_c\,\tfrac{1}{1+c(s)}
    + w_\kappa\,\kappa(s)
    + w_r\,\bigl(1-\max(\pi(s),h(s))\bigr) \nonumber\\
  &\quad
    - w_u\,\bigl(1-\text{conf}(s)\bigr),
  \label{eq:score}
\end{align}
where the latency component normalizes against the request ceiling:
\begin{equation}
  \ell\text{-comp}(s,q) =
    \begin{cases}
      1 & \text{if } \ell_{\max} = \infty, \\
      \max\!\bigl(0,\; 1 - \ell(s)/\ell_{\max}\bigr) & \text{otherwise.}
    \end{cases}
\end{equation}
Confidence $\text{conf}(s)$ is a weighted average of evidence confidence
values, where each item's confidence is
$n_{\text{eff}}/(n_{\text{eff}}+20)$ for effective count
$n_{\text{eff}} = n + \alpha_0 + \beta_0$.

\subsection{Pareto Annotation}

A source $a$ \emph{dominates} $b$ if $a$ is at least as good as $b$ on
every score component (trust, latency, freshness, cost, coverage, risk,
schema) and strictly better on at least one.  All non-dominated sources
are marked Pareto-optimal in the plan, enabling downstream policy
inspection and ablation without re-running the optimizer.

\subsection{Explanation and Audit}

Each \texttt{QueryPlan} ships a structured explanation listing the
top-$k$ sources with their score, component breakdown, Pareto status,
and top contributing dimensions.  Rejected sources are listed with the
violated constraints.  The \texttt{certify\_plan} API wraps the plan in
a stable, deterministically hashed certificate for deployment audit logs
and regression tests.

%------------------------------------------------------------------------
\section{Policy Profiles}
\label{sec:policies}

Named profiles make experiments repeatable without rewriting query
requests.  Table~\ref{tab:policies} shows the five profiles implemented
in the artifact.

\begin{table}[t]
\caption{Named policy profiles.  $w_\tau$, $w_\ell$, $w_c$, $w_r$:
  weights for trust, latency, cost, risk.
  $\tau_{\min}$: trust floor.  $\varepsilon$: risk tolerance.}
\label{tab:policies}
\centering
\begin{tabular}{lrrrrrr}
\toprule
Policy & $w_\tau$ & $w_\ell$ & $w_c$ & $w_r$ & $\tau_{\min}$ & $\varepsilon$\\
\midrule
high-assurance   & 0.56 & 0.08 & 0.04 & 0.13 & 0.72 & 0.05 \\
balanced         & 0.43 & 0.15 & 0.10 & 0.09 & 0.58 & 0.10 \\
latency-critical & 0.22 & 0.50 & 0.08 & 0.08 & 0.45 & 0.20 \\
cost-efficient   & 0.12 & 0.18 & 0.52 & 0.07 & 0.00 & 0.25 \\
trust-only       & 0.90 & 0.00 & 0.00 & 0.05 & 0.00 & 0.05 \\
\bottomrule
\end{tabular}
\end{table}

The profiles span the design space from maximum trust assurance to
pure cost minimization, enabling ablation studies that isolate the
contribution of each dimension.

%------------------------------------------------------------------------
\section{Extensions}
\label{sec:extensions}

\subsection{Portfolio Planning}

For high-stakes queries a single source may be insufficient.
\texttt{TrustAwarePortfolioPlanner} selects a small redundant set of
sources under a \texttt{PortfolioBudget} that specifies minimum and
maximum source counts, total latency, total cost, and a minimum
portfolio trust lower bound.  Portfolio trust uses a redundancy-aware
success model:
\begin{equation}
  \text{PTrust}(\mathbf{s}) =
    1 - \prod_{s \in \mathbf{s}}\bigl(1-\text{TLB}(s)\bigr)
    \cdot \delta^{|\mathbf{s}|-1},
  \label{eq:portfolio}
\end{equation}
where $\delta \in [0,1]$ is a configurable correlation penalty (default
$0.35$).  When $\delta = 1$ sources are fully independent; when
$\delta = 0$ adding sources provides no benefit.  This lets the artifact
study triangulated retrieval without pretending that evidence streams are
perfectly independent.

\subsection{Multi-Stage Pipeline Planning}

AI-native queries are often workflows rather than single-operator plans.
\texttt{TrustAwarePipelinePlanner} decomposes a query into stages
(retrieve, ground, generate, verify, policy audit) and applies a shared
policy to each stage.  The resulting \texttt{PipelinePlan} reports
end-to-end latency, total cost, completeness, and the minimum trust
lower bound across all stages---the \emph{pipeline trust floor}.

\subsection{Online Trust Updates}

\texttt{TrustLedger} maintains a log of \texttt{FeedbackEvent} records
(accepted / rejected) and folds new observations back into source
evidence streams.  This closes the feedback loop needed by AI-native
systems whose retrievers, tools, and model operators change behavior over
time.  The update is a simple append to the beta-binomial sufficient
statistics, preserving conjugacy and keeping the optimizer API unchanged.

\subsection{Trust Certificates}

\texttt{certify\_plan(plan, request)} produces a stable, deterministically
hashed audit record.  The certificate contains the selected source, ranked
alternatives, evidence summaries, hard constraints, Pareto status, and
rejection reasons.  \texttt{certify\_portfolio} applies the same idea to
redundant source sets.  Certificates are suitable for paper artifacts,
deployment audit logs, and regression tests.

%------------------------------------------------------------------------
\section{Evaluation}
\label{sec:eval}

\subsection{Experimental Setup}

All experiments run on a dependency-free Python implementation with
deterministic random seeds.  Synthetic experiments use 32 sources
generated from beta-distributed latent quality ($\text{Beta}(8,3)$),
with latency uniformly in $[20, 450]$\,ms, cost in $[0.02, 2.50]$, and
two evidence streams (\texttt{answer\_acceptance},
\texttt{policy\_compliance}) with 12--220 observations each.  The
workload contains 64 query requests cycling through six capability
combinations with randomly sampled latency, freshness, trust, and
privacy constraints.

\subsection{Synthetic Benchmark Ablation}
\label{sec:synth}

Table~\ref{tab:ablation} reports the four-policy ablation at seed 7.
The trust-aware policy selects the highest calibrated trust lower bound
(0.710) among policies that accept more than half the workload, at the
price of a modest latency increase over cost-first and fast-first.
Cost-first maximizes throughput (60/64 accepted) and minimizes cost
(0.651) but drops trust by 11.4 percentage points relative to the
trust-aware policy.  Fast-first achieves the lowest latency (85.0\,ms)
but incurs the highest cost (1.321).  Trust-only, while achieving the
highest absolute trust (0.737), accepts the fewest queries (51) because
its strict $\varepsilon = 0.05$ raises the trust lower bound threshold
for many sources.

\begin{table}[t]
\caption{Synthetic benchmark ablation (seed 7, 32 sources, 64 queries).
  \emph{acc.}: queries with at least one admissible source.
  \emph{score}: average weighted utility.
  \emph{TLB}: average calibrated trust lower bound.}
\label{tab:ablation}
\centering
\begin{tabular}{lrrrrr}
\toprule
Policy & Acc. & Score & TLB & Latency\,(ms) & Cost\\
\midrule
trust-aware & 53 & 0.733 & \textbf{0.710} & 112.8 & 0.829 \\
cost-first  & 60 & 0.693 & 0.629          & \textbf{79.2} & \textbf{0.651} \\
fast-first  & 60 & \textbf{0.808} & 0.689 &  85.0 & 1.321 \\
trust-only  & 51 & 0.733 & 0.737          & 168.6 & 1.050 \\
\bottomrule
\end{tabular}
\end{table}

\subsection{Multi-Seed Policy Study}
\label{sec:study}

Table~\ref{tab:study} reports average results over five random seeds,
confirming that per-seed observations are not artifacts.
High-assurance achieves the highest trust lower bound (0.815) at the
cost of accepting only 46.2 queries on average---its strict trust floor
($\tau_{\min} = 0.72$) and tight $\varepsilon = 0.05$ exclude sources
whose evidence streams are sparse.  Latency-critical achieves the lowest
average latency (108.2\,ms) and cost-efficient the lowest average cost
(0.665).  Balanced occupies the middle ground with 80.8 accepted queries,
a trust lower bound of 0.767, and the second-lowest latency.  These
results confirm that the named policies faithfully reflect their design
intent across seeds.

\begin{table}[t]
\caption{Multi-seed policy study (5 seeds, 32 sources, 64 queries/seed).
  All values are seed averages.  \emph{acc.}: average accepted queries.}
\label{tab:study}
\centering
\begin{tabular}{lrrrrr}
\toprule
Policy & Acc. & Score & TLB & Latency\,(ms) & Cost\\
\midrule
high-assurance   & 46.2 & 0.765          & \textbf{0.815} & 218.9 & 1.056\\
balanced         & 80.8 & 0.734          & 0.767          & 131.0 & 0.955\\
latency-critical & 83.2 & 0.717          & 0.740          & \textbf{108.2} & 1.086\\
cost-efficient   & 83.6 & 0.679          & 0.734          & 149.9 & \textbf{0.665}\\
trust-only       & 82.0 & \textbf{0.780} & 0.789          & 173.5 & 1.165\\
\bottomrule
\end{tabular}
\end{table}

\subsection{Real-Data Study: Wisconsin Diagnostic Breast Cancer}
\label{sec:wdbc}

\paragraph{Dataset.}
We use the UCI Wisconsin Diagnostic Breast Cancer (WDBC)
dataset~\cite{street1993nuclear} (569 instances, 30 features, 212
malignant / 357 benign).  Five feature-family \emph{diagnostic sources}
serve as a proxy for heterogeneous AI-native operators:
\texttt{compact-morphology} (4 mean-morphology features),
\texttt{nuclear-shape-risk} (4 shape-risk features),
\texttt{worst-case-specialist} (4 worst-case features),
\texttt{texture-screen} (3 texture features), and
\texttt{full-cytology-panel} (all 30 features).

\paragraph{Methodology.}
Each source is evaluated with deterministic stratified 5-fold
cross-validation.  A threshold linear model is fit on training folds
and evaluated on the test fold.  Aggregate accuracy, malignant
sensitivity, and benign specificity become the three evidence streams
fed to the trust optimizer.  Uncertainty is reported via Wilson score
intervals for binomial metrics and bootstrap 95\% CIs over fold-level
balanced accuracy (1\,000 resamples).

\paragraph{Source Reliability.}
Table~\ref{tab:wdbc_sources} reports reliability statistics.
\texttt{worst-case-specialist} achieves the best balanced accuracy
($92.3\%$, CI $[91.3, 93.5]$) with the lowest false-negative rate
(0.080), meaning only 8.0\% of malignant cases are missed.
\texttt{full-cytology-panel} follows at $88.5\%$ but requires all 30
features, entailing higher latency (86\,ms) and cost (0.18).
\texttt{texture-screen}, despite high sensitivity ($84.4\%$), has poor
specificity ($50.4\%$), creating an unacceptable false-positive burden.

\begin{table}[t]
\caption{WDBC source reliability (5-fold CV, Wilson 95\% intervals).
  \emph{bal\,acc}: balanced accuracy bootstrap 95\% CI.
  \emph{FNR}: false-negative rate.}
\label{tab:wdbc_sources}
\centering
\small
\begin{tabular}{lrccr}
\toprule
Source & Feat. & Bal\,acc CI\,(\%) & Sens CI\,(\%) & FNR\\
\midrule
worst-case-specialist  & 4  & $92.3\;[91.3,93.5]$ & $92.0\;[87.5,94.9]$ & 0.080\\
full-cytology-panel    & 30 & $88.5\;[87.1,90.1]$ & $87.3\;[82.1,91.1]$ & 0.127\\
compact-morphology     & 4  & $86.6\;[85.6,87.7]$ & $87.3\;[82.1,91.1]$ & 0.127\\
nuclear-shape-risk     & 4  & $84.2\;[82.9,85.4]$ & $82.1\;[76.4,86.7]$ & 0.179\\
texture-screen         & 3  & $67.4\;[65.5,69.5]$ & $84.4\;[78.9,88.7]$ & 0.156\\
\bottomrule
\end{tabular}
\end{table}

\paragraph{Policy Decisions.}
Table~\ref{tab:wdbc_decisions} shows which source each policy selects.
The high-assurance and trust-only policies both select
\texttt{worst-case-specialist} (score 0.883 and 0.900 respectively),
reflecting its dominant balanced accuracy and lowest false-negative rate.
The balanced, latency-critical, and cost-efficient policies select
\texttt{compact-morphology}, which offers competitive sensitivity (87.3\%)
with the lowest latency (22\,ms) and cost (0.04).  This is the correct
engineering trade-off: in a latency-sensitive or cost-constrained
deployment, a source 5.7 percentage points behind on balanced accuracy
but 4$\times$ faster and 2.5$\times$ cheaper is a rational choice.

The results illustrate the core value of the optimizer: policy
\emph{intent} is transparently mapped to source \emph{selection} through
auditable weights and constraints, rather than being hidden in
application-level if-statements.

\begin{table}[t]
\caption{WDBC policy decisions.  \emph{TLB}: calibrated trust lower
  bound.  \emph{Sens}/\emph{Spec}: malignant sensitivity/benign specificity.}
\label{tab:wdbc_decisions}
\centering
\begin{tabular}{llrrrrr}
\toprule
Policy & Selected & Score & TLB & Sens & Spec & FNR\\
\midrule
high-assurance   & worst-case-specialist & 0.883 & \textbf{0.898} & \textbf{0.920} & \textbf{0.927} & \textbf{0.080}\\
balanced         & compact-morphology    & 0.879 & 0.838 & 0.873 & 0.860 & 0.127\\
latency-critical & compact-morphology    & 0.840 & 0.838 & 0.873 & 0.860 & 0.127\\
cost-efficient   & compact-morphology    & \textbf{0.911} & 0.838 & 0.873 & 0.860 & 0.127\\
trust-only       & worst-case-specialist & 0.900 & \textbf{0.898} & \textbf{0.920} & \textbf{0.927} & \textbf{0.080}\\
\bottomrule
\end{tabular}
\end{table}

\subsection{Discussion}

The evaluation validates four claims:

\smallskip\noindent\textbf{C1: Trust lower bounds are meaningful.}
Sources with sparse evidence receive lower bounds conservatively below
their declared score (Eq.~\ref{eq:tlb}).  \texttt{worst-case-specialist}
earns TLB~0.898 from 212 positive malignant observations; a source with
only 12 observations cannot reach the same floor regardless of its
observed rate.

\smallskip\noindent\textbf{C2: Hard constraints prevent trust leakage.}
In the WDBC study, \texttt{texture-screen} is rejected by the
high-assurance policy because its hallucination risk (proxy:
false-negative rate 0.156) exceeds the ceiling.  Its high sensitivity
does not compensate.

\smallskip\noindent\textbf{C3: Policy profiles produce faithful trade-offs.}
Across five seeds and five policies, the winners for each metric are
exactly the policies designed to win it (Table~\ref{tab:study}).

\smallskip\noindent\textbf{C4: The framework generalizes across domains.}
The same optimizer handles random synthetic catalogs and the WDBC
real-data study without code changes, confirming that the model is not
over-fit to a specific domain.

%------------------------------------------------------------------------
\section{Implementation Notes}
\label{sec:impl}

The artifact is a pure-Python package with no runtime dependencies.
Statistical primitives (Wilson intervals, bootstrap CIs, beta posteriors)
are implemented directly.  The package exposes a typed public API
compatible with Python\,3.9+ and enforces immutability throughout via
frozen dataclasses.

Three scoring strategies are available: \texttt{linear} (default, as
described in \S\ref{sec:optimizer}), \texttt{topsis} (Technique for
Order Preference by Similarity to an Ideal Solution), and
\texttt{bayesian-ucb} (Upper Confidence Bound).  All strategies share
the same constraint-filtering front end; only the ranking function
differs.

The visualization module generates publication-quality dependency-free
SVG figures covering the policy trade-off atlas, Pareto skyline, score
waterfall, multi-stage AI pipeline diagram, calibration lens, and
real-data confidence interval summary.

%------------------------------------------------------------------------
\section{Conclusion}
\label{sec:conclusion}

We have presented trust-aware query optimization, an extension of
classical cost-based planning that treats calibrated source trust as a
first-class objective.  The key ideas are:
\begin{enumerate}
  \item \emph{Calibrated lower bounds} from beta-binomial posteriors
    prevent optimistic trust scores from winning under sparse evidence.
  \item \emph{Hard policy gates} keep constraints auditable rather than
    hiding them in scalar compensation.
  \item \emph{Explainable plan objects} and \emph{portable certificates}
    turn every optimizer decision into an audit record.
  \item \emph{Named policy profiles} make experiments and deployments
    reproducible.
  \item \emph{Portfolio, pipeline, and feedback extensions} lift
    single-source planning into the full AI query workflow lifecycle.
\end{enumerate}

Experiments on a 32-source synthetic benchmark and the WDBC real-data
study confirm that the trust-aware policy achieves the best calibrated
trust lower bound among utility-maximising policies and that hard
constraints prevent trust leakage even when other metrics appear
attractive.  The framework applies unchanged to both synthetic catalogs
and medical-diagnostic feature-family sources, demonstrating
generalization across domains.

Future work includes extending the beta-binomial model to multi-label
trust (e.g., separate dimensions for factual accuracy, copyright
compliance, and safety); integrating with real RAG pipeline telemetry;
and studying the interaction between portfolio redundancy and correlation
structure in practical deployments.

%------------------------------------------------------------------------
\bibliographystyle{splncs04}
\begin{thebibliography}{99}

\bibitem{lewis2020rag}
Lewis, P., et al.:
Retrieval-augmented generation for knowledge-intensive NLP tasks.
In: Advances in Neural Information Processing Systems. vol.~33,
pp.~9459--9474 (2020)

\bibitem{borgeaud2022retro}
Borgeaud, S., et al.:
Improving language models by retrieving from trillions of tokens.
In: International Conference on Machine Learning,
pp.~2206--2240. PMLR (2022)

\bibitem{selinger1979access}
Selinger, P.G., et al.:
Access path selection in a relational database management system.
In: Proceedings of ACM SIGMOD, pp.~23--34 (1979)

\bibitem{blaze1996decentralized}
Blaze, M., Feigenbaum, J., Lacy, J.:
Decentralized trust management.
In: Proceedings of IEEE Symposium on Security and Privacy,
pp.~164--173 (1996)

\bibitem{buneman2001provenance}
Buneman, P., Khanna, S., Tan, W.:
Why and where: A characterization of data provenance.
In: International Conference on Database Theory,
pp.~316--330. Springer (2001)

\bibitem{batini2009methodologies}
Batini, C., Scannapieco, M.:
Methodologies for data quality assessment and improvement.
ACM Computing Surveys~\textbf{41}(3), 1--52 (2009)

\bibitem{levy1996querying}
Levy, A.Y., Rajaraman, A., Ordille, J.J.:
Querying heterogeneous information sources using source descriptions.
In: Proceedings of Very Large Data Bases, pp.~251--262 (1996)

\bibitem{antova2008maybms}
Antova, L., Koch, C., Olteanu, D.:
From complete to incomplete information and back.
In: Proceedings of ACM SIGMOD, pp.~713--726 (2008)

\bibitem{widom2005trio}
Widom, J.:
Trio: A system for data, uncertainty, and lineage.
In: Proceedings of Very Large Data Bases, pp.~1497--1500 (2005)

\bibitem{kamvar2003eigentrust}
Kamvar, S.D., Schlosser, M.T., Garcia-Molina, H.:
The EigenTrust algorithm for reputation management in P2P networks.
In: Proceedings of the 12th WWW Conference, pp.~640--651 (2003)

\bibitem{winsborough2002trust}
Winsborough, W.H., Li, N.:
Towards practical automated trust negotiation.
In: Proceedings of the 3rd Workshop on Policies for Distributed
Systems and Networks, pp.~92--103 (2002)

\bibitem{re2007probabilistic}
R\'{e}, C., Dalvi, N., Suciu, D.:
Efficient top-$k$ query evaluation on probabilistic data.
In: Proceedings of IEEE ICDE, pp.~886--895 (2007)

\bibitem{ioannidis1990randomized}
Ioannidis, Y., Kang, Y.C.:
Randomized algorithms for optimizing large join queries.
In: Proceedings of ACM SIGMOD, pp.~312--321 (1990)

\bibitem{mohan1992aries}
Mohan, C., et al.:
ARIES: A transaction recovery method supporting fine-granularity locking
and partial rollbacks using write-ahead logging.
ACM Transactions on Database Systems~\textbf{17}(1), 94--162 (1992)

\bibitem{patel2024lotus}
Patel, L., et al.:
LOTUS: Enabling semantic queries with LLMs over tables of unstructured
and structured data.
arXiv preprint arXiv:2407.11418 (2024)

\bibitem{shi2024survey}
Shi, P., et al.:
A survey of NL2SQL with large language models: Where are we, and where
are we going?
arXiv preprint arXiv:2408.05109 (2024)

\bibitem{street1993nuclear}
Street, W.N., Wolberg, W.H., Mangasarian, O.L.:
Nuclear feature extraction for breast tumor diagnosis.
In: IS\&T/SPIE International Symposium on Electronic Imaging,
vol.~1905, pp.~861--870 (1993)

\end{thebibliography}

\end{document}
